%%%%%%%%%%%%%%%%%%%%%%%%%%%%%%%%%%%%%%%%%%%%%%%%%%%%%%%%%%%%%%%%%%%%%%%%%%
%
%                         华南理工大学文献综述模板
%
%   本文件为模板示例，请根据个人需求进行相应修改。
%
%   1. 请使用 VS Code“打开文件夹”打开“文献综述”目录作为工作空间。
%      本目录的 .vscode/settings.json 会自动生效，无需复制或修改配置。
%      Ctrl+S 保存时默认快速编译；首次使用或更新文献后请手动完整编译。
%   2. 本文件填写封面信息、摘要及正文；文献条目在 main.bib 中维护。
%      封面、说明页和排版由 scutreview.cls 生成，通常无需修改类文件。
%   3. degree 控制封面的硕士/博士方框；留空 {} 时不勾选。
%      instructions = true 保留“说明”页，改为 false 可去掉该页。
%   4. 字体位于上一级 fonts/，校名图片位于上一级 public/。
%      编译结果保存到本目录 build/；请保留项目原有目录结构。
%
%%%%%%%%%%%%%%%%%%%%%%%%%%%%%%%%%%%%%%%%%%%%%%%%%%%%%%%%%%%%%%%%%%%%%%%%%%

\documentclass{scutreview}

\addbibresource{main.bib} % 参考文献，在 main.bib 中收录

%%%%%%%%%%%%%%%%%%%%%%%%%%%%%%%%%%%%%%%%%%%%%%
% 封面、正文题目及勾选项
%%%%%%%%%%%%%%%%%%%%%%%%%%%%%%%%%%%%%%%%%%%%%%
\reviewsetup{
  degree       = {},     % master：硕士；doctor：博士；填写的时候带大括号，例如 {master}
  author       = {姓名}, % 研究生姓名，同时用于正文署名
  student-id   = {学号},
  department   = {院（系）名称},
  major        = {专业名称},
  supervisor   = {指导教师姓名},
  date         = {二〇二六年四月},
  review-title = {请填写文献综述题目},           % 正文标题
  keywords     = {关键词一；关键词二；关键词三}, % 关键词
  instructions = true                            % true：保留说明页；false：去掉说明页
}

\begin{document}

\maketitle

%%%%%%%%%%%%%%%%%%%%%%%%%%%%%%%%%%%%%%%%%%%%%%
% 摘要
%%%%%%%%%%%%%%%%%%%%%%%%%%%%%%%%%%%%%%%%%%%%%%
% 摘要不超过 500 字
\begin{reviewabstract}
请在此撰写中文摘要，概括综述主题、文献范围、主要研究脉络、核心争议与总体结论，字数不超过 500 字。
\end{reviewabstract}

%%%%%%%%%%%%%%%%%%%%%%%%%%%%%%%%%%%%%%%%%%%%%%
% 关键词
%%%%%%%%%%%%%%%%%%%%%%%%%%%%%%%%%%%%%%%%%%%%%%
% 定义在了 “封面、正文题目及勾选项” 的 \reviewsetup 中
% 在 keywords 字段填写即可，此处不需要修改！！！

%%%%%%%%%%%%%%%%%%%%%%%%%%%%%%%%%%%%%%%%%%%%%%
% 一、前言
%%%%%%%%%%%%%%%%%%%%%%%%%%%%%%%%%%%%%%%%%%%%%%
% 自动生成编号！！！后文同理
\section{前言}

请说明文献综述的写作目的与意义，介绍有关概念并提供必要的背景材料；描述课题的国内外研究现状、相关研究的演变与最新进展，归纳尚存问题、争论焦点和发展趋势，并交代本综述讨论的范围。

%%%%%%%%%%%%%%%%%%%%%%%%%%%%%%%%%%%%%%%%%%%%%%
% 二、正文
%%%%%%%%%%%%%%%%%%%%%%%%%%%%%%%%%%%%%%%%%%%%%%
\section{正文}

% （一）历史发展
\subsection{历史发展}

请梳理研究主题的发展阶段、代表性理论或技术路线，以及各阶段达到的研究水平\cite{sample1}。

% （二）现状分析
\subsection{现状分析}

请分别介绍国外与国内研究现状，从研究问题、方法、数据、评价指标和应用场景等方面比较研究差距，指出国内外研究空白及未来发展趋势\cite{sample2}。

% （三）趋向预测
\subsection{趋向预测}

请在纵向演进与横向比较的基础上，评价当前研究水平、现存问题和不同意见，并提出有依据的展望。

%%%%%%%%%%%%%%%%%%%%%%%%%%%%%%%%%%%%%%%%%%%%%%
% 三、总结与展望
%%%%%%%%%%%%%%%%%%%%%%%%%%%%%%%%%%%%%%%%%%%%%%
\section{总结与展望}

请高度概括综述主题及主要结论，指出既有研究的不足和尚未解决的问题，说明现有成果的意义与价值，并明确未来值得继续研究的问题和方向。

%%%%%%%%%%%%%%%%%%%%%%%%%%%%%%%%%%%%%%%%%%%%%%
% 参考文献列表
%%%%%%%%%%%%%%%%%%%%%%%%%%%%%%%%%%%%%%%%%%%%%%
% 新增、删除或修改引用后，请手动选择“完整编译”以运行 Biber 并更新列表。
\printbibliography
\end{document}
